% Working exact free-lattice input for the maps-owned OS45 invariant problem.
% No complete torsion fundamental-class invariant or smooth classification is asserted.
\subsubsection*{The primitive component lattice and cubic in the sixteen-component OS45 example}
We first record a saturation calculation which uses the actual component
meridians. Let $X$ be simply connected, let $N$ be a single rank-two
Mumford cusp with $D$ components, and put $E=X\setminus\operatorname{int}N$.
Assume $b_2(X)=D-1$, $H_3(N;\mathbb Z)$ is free, and
\[
 H_1(E;\mathbb Z)=\mathbb Z\oplus A,\qquad
                       \mu_i=(1,a_i),                      \tag{D.PC1}
\]
where $A$ is finite and $\mu_i$ is the positively oriented meridian of
component $i$. The component classes $x_i$ then generate a sublattice
\[
 L=\mathbb Z^D/\mathbb Z(1,\ldots,1)
                       \subset H^2(X;\mathbb Z)
\]
whose full primitive extension, inside $L\otimes\mathbb Q$, is
\[
 H^2(X;\mathbb Z)
   =L+\left\{\sum_i\widetilde{\chi(a_i)}x_i:
                         \chi\in\operatorname{Hom}(A,\mathbb Q/\mathbb Z)
        \right\}.                                        \tag{D.PC2}
\]
A tilde means any rational lift of a value in $\mathbb Q/\mathbb Z$;
changing it adds an element of $L$. In particular the component
lattice has index $|A|$.

\begin{proof}
Poincare--Lefschetz duality identifies
$H^2(X,E;\mathbb Z)=H_4(N;\mathbb Z)=\mathbb Z^D$,
with its positive component Thom basis, and
$H^3(X,E;\mathbb Z)=H_3(N;\mathbb Z)$ is free. The connecting
map from $H^1(E;\mathbb Z)$ evaluates on the meridians (D.PC1),
so its image is precisely $\mathbb Z(1,\ldots,1)$.
The relative exact sequence therefore injects $L$ into $H^2(X)$.
Since $X$ is simply connected, the latter group is free; its rank is
$D-1$ by hypothesis. Its quotient by $L$ is finite and embeds in
$H^2(E)$, whose torsion is
$\operatorname{Ext}(H_1(E),\mathbb Z)=\operatorname{Hom}(A,\mathbb Q/\mathbb Z)$.
All this torsion lifts to $H^2(X)$ because the next relative group is
free. Thus the quotient has exactly that order.

Here is also its precise embedding, rather than only its index.
For a character $\chi$ of order $m$, extend it by zero on the free
summand of $H_1(E)$. Regard it as a class in
$H^1(E;\mathbb Z/m)$ by multiplying its values by $m$.
Its relative coboundary has component coordinates
$m\widetilde{\chi(a_i)}\pmod m$. Choose the displayed integers
as an integral relative lift. Its image in $H^2(X;\mathbb Z/m)$
is zero by exactness. Hence its integral image is divisible by $m$
in $H^2(X;\mathbb Z)$, giving the class in (D.PC2).
The coefficient exact sequences identify its exterior restriction
with the corresponding torsion Bockstein, up to the harmless common
connecting-map sign. Since these characters give the full torsion
of $H^2(E)$, they give the full primitive extension. This proves
(D.PC2), including every finite denominator.
\end{proof}

Now take the exact marked data
\[
 \mathrm{OS45/default},\quad P=2G,\quad Q=8G,\quad
 (w_0,w_1,w_2,w_3,w_4)=(2,0,0,0,0),\quad e=1,              \tag{D.PC3}
\]
with all multiplicative fibres filled by their marked Mumford caps.
D.2.46 gives
\[
 H_2(X)=\mathbb Z^{15}\oplus\mathbb Z/8,\qquad
 H_3(X)=\mathbb Z/8,\qquad \pi_1(X)=1.
                                                               \tag{D.PC4}
\]
The elliptic monodromy at $I_8$ is
$\left(\begin{smallmatrix}1&8\\0&1\end{smallmatrix}\right)$
and its section cocycle is $(3,0)$. Formula (D.GI2) therefore gives
\[
 A_0=\begin{pmatrix}
 1&8&0&24\\0&1&0&0\\0&-6&1&-16\\0&0&0&1
 \end{pmatrix},\qquad
 L_{\rm per}=\left\langle(8,-6),(24,-16)\right\rangle
                 =8\mathbb Z\oplus2\mathbb Z.             \tag{D.PC5}
\]
The short basis is $(e_1,\delta)$ and the long basis is $(e_2,\gamma)$.
This is an equality of sublattices, not merely a comparison of
indices. The determinant is sixteen and the Smith factors are two
and eight. Label the normalized components by
\[
             (r,s)\in\mathbb Z/8\times\mathbb Z/2,
             \qquad i=2r+s.                              \tag{D.PC6}
\]

The actual cusp complement satisfies
\[
 H_1(E)=\mathbb Z\oplus\mathbb Z/2,
                      \qquad \mu_{r,s}=(1,s\bmod2).      \tag{D.PC7}
\]
To verify the marking, retain all peripheral generators and the
ordinary clutch in the exact degree-one presentation. Its Smith
form has seven unit factors, one factor two, and one free generator.
The torsion character in the original nine coordinates is
\[
                  (0,-2,1,2,0,0,2,-2,2)\pmod2,
                                                               \tag{D.PC8}
\]
which is precisely $\delta^*$ modulo two. The free character
evaluates to one on the cusp meridian. A height-one ray at
$(r,s)$ has literal meridian $\theta+r e_1+s\delta$; it
therefore gives exactly (D.PC7). The full original presentation,
Smith change of basis and all sixteen meridian images are saved
in the executable below.

Applying (D.PC2), there is a unique integral class $y$ with
\[
                    2y=\sum_{r=0}^{7}x_{r,1}.             \tag{D.PC9}
\]
The component lattice has exact index two. In particular it would
be incorrect to declare any fifteen of the component classes a
primitive basis. A convenient actual basis is
\[
 x_0,x_1,\ldots,x_{13},y,\qquad
 x_{14}=-2y-\sum_{\substack{i<14\\i\ \mathrm{even}}}x_i,
 \qquad
 x_{15}=2y-\sum_{\substack{i<14\\i\ \mathrm{odd}}}x_i.
                                                               \tag{D.PC10}
\]
The integral matrix of the sixteen component classes in this basis
has Smith factors $1^{14},2$ and kernel generated by the all-one
vector. These are the asserted primitive markings.

Here is the complete cubic on that basis. For component classes,
D.2.47 specializes to the following symmetric integral trilinear
form. With distinct vertices as appropriate,
\[
 \begin{aligned}
 T(x_i,x_i,x_i)&=6,\\
 T(x_i,x_i,x_j)&=
 \begin{cases}
 -2,&i=(r,s),\ j=(r,1-s),\\
 -1,&i=(r,s),\ j=(r\pm1,t),\quad t=0,1,\\
 0,&\text{otherwise},
 \end{cases}\\
 T(x_i,x_j,x_k)&=
 \begin{cases}
 1,&\{i,j,k\}\text{ is three of the four vertices in}\\[-2pt]
   &\hspace{12pt}\{(r,0),(r,1),(r+1,0),(r+1,1)\}
                   \text{ for some }r,\\
 0,&\text{otherwise}.
 \end{cases}
 \end{aligned}                                             \tag{D.PC11}
\]
All horizontal indices are cyclic modulo eight. Repeated geometric
edges and triangle cells are retained in this formula.
The remaining primitive entries are
\[
 \begin{aligned}
 T(y,x_{r,s},x_{r',s'})
 &=(-1)^s\left(
       \mathbf1_{s'=s,\ r'=r+1}
      +\mathbf1_{s'=s,\ r'=r-1}
      -2\mathbf1_{(r',s')=(r,s)}\right),\\
 T(y,y,x_i)&=0,\qquad T(y,y,y)=0.
 \end{aligned}                                             \tag{D.PC12}
\]
These formulas determine every entry on the primitive basis
(D.PC10), and every resulting entry is integral.

For a geometric verification of the half-class entries, on a
normalized $dP_6$ take the cyclic boundary classes
\[
 E_1,\ H-E_1-E_2,\ E_2,\ H-E_2-E_3,\ E_3,\ H-E_1-E_3
                                                               \tag{D.PC13}
\]
in the ray order of D.2.47. The four nonhorizontal boundary classes
sum to $2f$, where $f=H-E_1$ and $f^2=0$. Restricting (D.PC9)
therefore gives
\[
                         y|_{P_{r,s}}=(-1)^s f.            \tag{D.PC14}
\]
This division is valid in the torsion-free group $H^2(dP_6;\mathbb Z)$.
The fibre $f$ has degree one on each horizontal boundary curve
and zero on the other four. The Cartier restriction formula in
D.2.47 now gives the first line of (D.PC12). Squaring (D.PC14)
gives $y^2x_i=0$; multiplying (D.PC9) by $y^2$ gives $y^3=0$.
Thus no rational cubic has been silently substituted for the
integral form.

Finally $c_1(X)=0$ and both $c_2$ and $p_1$ pair to zero with
every class in the primitive free basis. For $c_1$, the normalized
component tangent calculation is zero and the component restriction
is rationally injective by the connected toric graph Laplacian;
$H^2(X)$ is torsion-free, so this proves integral vanishing.
For $c_2,p_1$, D.2.47 gives zero on every component. The component
classes have finite index, and the evaluations are integers, so they
also vanish on $y$. These statements concern the free characteristic
pairings; they do not determine the torsion class of $c_2$, the
refined class, or the full finite-coefficient fundamental-class
invariant.

\path{code/fourfolds_OS45_primitive_cubic.py}
reconstructs the actual complement before using its torsion
character, checks all sixteen meridians, computes the exact
index-two primitive lattice, and verifies all $15^3=3375$ primitive
triples and all $16^3=4096$ recovered component triples. The
normalization restriction (D.PC14), the actual period basis, and
all sparse integral entries are retained in its JSON output.
The preliminary component-only calculation is retained separately
in \path{code/fourfolds_OS45_component_ring.py}.
